%%
%% MScheme: Comprehensive User Manual
%%
\documentclass[12pt]{report}
\usepackage[margin=1in]{geometry}
\usepackage{longtable}
\usepackage{booktabs}
\usepackage{listings}
\usepackage{xcolor}
\usepackage{hyperref}
\usepackage{textcomp}

\lstdefinelanguage{Scheme}{
  morekeywords={define,lambda,if,cond,else,let,let*,letrec,begin,
    set!,quote,quasiquote,unquote,unquote-splicing,and,or,not,
    case,do,delay,force,map,for-each,apply,call-with-current-continuation,
    call/cc,require-modules,load,macro,define-global-symbol,
    import,module,library,program,SchemeStubs,Smodule,ExportSchemeStubs,
    importSchemeStubs},
  sensitive=true,
  morecomment=[l]{;},
  morecomment=[l]{\%},
  morestring=[b]",
  literate={lambda}{$\lambda$}1,
}

\lstdefinelanguage{Modula3}{
  morekeywords={MODULE,INTERFACE,IMPORT,FROM,PROCEDURE,BEGIN,END,
    VAR,CONST,TYPE,RECORD,OBJECT,METHODS,OVERRIDES,RETURN,
    IF,THEN,ELSE,ELSIF,WHILE,DO,FOR,TO,BY,LOOP,EXIT,REPEAT,UNTIL,
    CASE,OF,WITH,TRY,EXCEPT,FINALLY,RAISE,RAISES,
    NEW,ARRAY,REF,REFANY,SET,BOOLEAN,INTEGER,REAL,LONGREAL,CHAR,TEXT,
    TRUE,FALSE,NIL,BRANDED,REVEAL,EVAL,NOT,AND,OR,IN,
    ISTYPE,NARROW,TYPECODE},
  sensitive=true,
  morecomment=[n]{(*}{*)},
  morestring=[b]",
}

\lstset{
  basicstyle=\ttfamily\small,
  breaklines=true,
  frame=single,
  backgroundcolor=\color{gray!5},
  keywordstyle=\bfseries\color{blue!70!black},
  commentstyle=\itshape\color{green!50!black},
  stringstyle=\color{red!60!black},
  showstringspaces=false,
  columns=flexible,
  tabsize=2,
  xleftmargin=2em,
  framexleftmargin=1.5em,
}

\title{\bfseries\sffamily\fontsize{24.88}{30}\selectfont
  MScheme User Manual}
\author{Mika Nystr\"om \\ {\tt mika@alum.mit.edu}}
\date{March 2026}

\begin{document}
\maketitle
\thispagestyle{empty}

\newpage
\tableofcontents
\newpage

%======================================================================
\chapter{Introduction}
%======================================================================

MScheme is a Scheme interpreter written in Modula-3, designed to serve
as an embedded scripting and glue language for Modula-3 applications.
It originated as a line-for-line translation of Peter Norvig's JScheme
(Java, circa~1998) into Modula-3, and has been substantially extended
with an automatic stub generator, a module system, and numerous
Modula-3 integration features.

MScheme follows the \emph{Revised$^4$ Report on the Algorithmic Language
Scheme (R4RS)}~\cite{r4rs} with extensions.  It is not a complete R5RS
implementation.  Notable limitations inherited from JScheme include:

\begin{itemize}
\item Only one numeric type: IEEE~754 double precision ({\tt LONGREAL}).
\item {\tt call-with-current-continuation} only unwinds the stack
  (no full continuations).
\item The Modula-3 call stack is used for Scheme control flow, limiting
  recursion depth to available stack space.
\end{itemize}

Despite these limitations, MScheme has been used to build substantial
production systems in financial trading, VLSI and circuit design (the Intel
asynchronous circuit tools), and scientific computing.

\section{Key Architectural Insight}

The central design feature of MScheme is that Modula-3's universal
traced reference type {\tt REFANY} serves directly as the Scheme value
type:

\begin{lstlisting}[language=Modula3]
(* SchemeObject.i3 *)
TYPE T = REFANY;
\end{lstlisting}

Because every Modula-3 heap-allocated object is already a valid {\tt REFANY},
it is already a valid Scheme value.  No wrapping, boxing, or marshaling
is needed to pass Modula-3 objects into Scheme or vice versa.  The Modula-3
garbage collector manages both Modula-3 and Scheme data uniformly.

This ``zero-impedance'' property distinguishes MScheme from
C-hosted embedding systems like Lua, Guile, Tcl, and CPython, where
every boundary crossing requires explicit conversion.  It is shared
only with the earlier DEC SRC ``Tinylisp'' system (Modula-2+, 1989)
and with JScheme and Kawa on the JVM.

\section{Components}

The MScheme system consists of several CM3 packages:

\begin{center}
\begin{tabular}{ll}
\toprule
\textbf{Package} & \textbf{Description} \\
\midrule
{\tt mscheme}         & Core Scheme interpreter \\
{\tt modula3scheme}   & Modula-3 integration layer \\
{\tt sstubgen}        & Stub generator tool and build-system templates \\
{\tt scheme-lib}      & Bundled Scheme library files \\
{\tt schemereadline}  & Interactive readline interface \\
{\tt schemesig}       & Signal handling \\
{\tt mscheme-interactive}   & Basic interactive interpreter \\
{\tt mscheme-interactive\_r} & Interactive interpreter with GNU Readline \\
{\tt schemecompiler}  & Ahead-of-time Scheme-to-Modula-3 compiler \\
\bottomrule
\end{tabular}
\end{center}


%======================================================================
\chapter{Getting Started}
%======================================================================

\section{Running the Interactive Interpreter}

The simplest way to use MScheme is to run one of the interactive
interpreters:

\begin{lstlisting}[language={}]
$ mscheme-interactive_r
>
\end{lstlisting}

This gives you a Scheme REPL (Read-Eval-Print Loop).  The prompt is
``{\tt >}''.  The result of the last evaluated expression is always
stored in the variable {\tt bang-bang} (interactive REPL only,
not available in loaded files).

\section{Environment Variables}

\begin{description}
\item[{\tt MSCHEMEPATH}] Colon-separated list of directories to search
  when {\tt load} is called.  Default is ``{\tt .}'' (current directory).
\item[{\tt NOMSCHEMETRACEBACKS}] If set (to any value), disables error
  tracebacks.
\end{description}

\section{Loading Files}

\begin{lstlisting}[language=Scheme]
(load "myfile.scm")
\end{lstlisting}

If a compiled version of the file is registered (see
Section~\ref{sec:registry}), it is used directly---no file is read.
Otherwise, files are searched along {\tt **scheme-load-path**}, which
is initialized from {\tt MSCHEMEPATH}.  If the file is not found on
disk, MScheme checks its bundled library (see
Chapter~\ref{ch:libraries}).
The load path can be extended at runtime:

\begin{lstlisting}[language=Scheme]
(set! **scheme-load-path**
      (cons "/my/extra/path" **scheme-load-path**))
\end{lstlisting}

Absolute paths and tilde-expansion are supported:

\begin{lstlisting}[language=Scheme]
(load "/usr/local/share/mylib.scm")
(load "~/scripts/util.scm")
\end{lstlisting}

\section{The Module System}\label{sec:modules}

MScheme provides a simple module system via {\tt require-modules}:

\begin{lstlisting}[language=Scheme]
(require-modules "basic-defs" "m3" "hashtable" "set" "display")
\end{lstlisting}

Each argument is either a bundled library name (without {\tt .scm}
extension) or a file path.  Modules are loaded at most once---repeated
{\tt require-modules} calls for an already-loaded module are no-ops.
This prevents loading loops.

\medskip\noindent\textbf{Important:}\ The {\tt require-modules} function is not available
by default in a bare {\tt Scheme.T}.  It is loaded when the string
{\tt "require"} is passed as the first file to the interpreter's
{\tt init()} method.  The {\tt SchemeM3.T} and the interactive
interpreters load it automatically.


%======================================================================
\chapter{Scheme Language Reference}
%======================================================================

MScheme implements R4RS with extensions.  This chapter documents
the differences from and additions to standard Scheme.

\section{Special Forms}

The following special forms are recognized by the evaluator:

\begin{center}
\begin{tabular}{ll}
\toprule
\textbf{Form} & \textbf{Description} \\
\midrule
{\tt (quote x)}                  & Standard quote \\
{\tt (if test then else)}        & Conditional \\
{\tt (cond clause \ldots)}       & Multi-way conditional (with {\tt else} and {\tt =>}) \\
{\tt (begin expr \ldots)}        & Sequential evaluation \\
{\tt (define var val)}           & Definition \\
{\tt (define (name args) body)}  & Procedure definition sugar \\
{\tt (set! var val)}             & Assignment \\
{\tt (lambda args body)}         & Closure creation \\
{\tt (macro args body)}          & Macro definition (non-hygienic) \\
{\tt (eval expr)}                & Evaluate in current environment \\
{\tt (current-environment)}      & Return current lexical environment \\
{\tt (unwind-protect e c h)}     & Exception handling \\
\bottomrule
\end{tabular}
\end{center}

\subsection{Macros}

MScheme uses the JScheme-style {\tt macro} form, which is a
non-hygienic, procedural macro:

\begin{lstlisting}[language=Scheme]
(define when
  (macro (test . body)
    `(if ,test (begin ,@body))))
\end{lstlisting}

The standard macros {\tt let}, {\tt let*}, {\tt letrec}, {\tt case},
{\tt do}, {\tt delay}, {\tt and}, {\tt or}, and {\tt quasiquote} are
all defined as macros in {\tt SchemePrimitives.i3}.  The {\tt time}
macro measures execution time:

\begin{lstlisting}[language=Scheme]
(time (fibonacci 30))
\end{lstlisting}

R5RS hygienic macros ({\tt syntax-rules}) are available via:

\begin{lstlisting}[language=Scheme]
(require-modules "mbe")
\end{lstlisting}

\subsection{{\tt unwind-protect}}

MScheme provides an {\tt unwind-protect} special form for exception handling:

\begin{lstlisting}[language=Scheme]
(unwind-protect
  protected-expression
  cleanup-expression     ; always runs (like "finally")
  error-expression)      ; runs only on exception
\end{lstlisting}


\section{Built-In Primitives}\label{sec:primitives}

\subsection{Arithmetic and Comparison}

\begin{lstlisting}[language=Scheme]
+  -  *  /  =  <  >  <=  >=
abs  floor  ceiling  round  truncate
max  min  quotient  remainder  div  modulo
sqrt  expt  exp  log
sin  cos  tan  asin  acos  atan
sinh  cosh  tanh  asinh  acosh  atanh
gcd  lcm
odd?  even?  zero?  positive?  negative?
number?  integer?  exact?  inexact?
exact->inexact  inexact->exact
number->string  string->number
\end{lstlisting}

All numbers are IEEE~754 double precision.  {\tt exact?}, {\tt
rational?}, and {\tt integer?} all test whether the number has no
fractional part.  {\tt complex?}, {\tt real?}, and {\tt number?} are
equivalent.  {\tt exact->inexact} and {\tt inexact->exact} are
identity functions.

{\tt atan} accepts one or two arguments: {\tt (atan $x$)} computes
$\arctan x$; {\tt (atan $y$ $x$)} computes $\arctan(y/x)$ with
correct quadrant (delegates to {\tt Math.atan2}).

Number literals may carry R4RS radix prefixes: {\tt \#b} (binary),
{\tt \#o} (octal), {\tt \#d} (decimal, default), and {\tt \#x}
(hexadecimal).  Exactness prefixes {\tt \#e} and {\tt \#i} are
accepted but are no-ops.  Combined prefixes such as {\tt \#e\#xFF} or
{\tt \#x\#e1F} are supported in any order.

\subsection{Pairs and Lists}

\begin{lstlisting}[language=Scheme]
cons  car  cdr  set-car!  set-cdr!
cadr  caddr  caar  cdar  cddr  ; ...all c[ad]{2,4}r
pair?  null?  list?
list  length  append  reverse
list-ref  list-tail
member  memq  memv  assoc  assq  assv
map  for-each  apply
_list*  ; (_list* 1 2 3) => (1 2 . 3)
\end{lstlisting}

\subsection{Booleans, Strings, Characters, Vectors, Symbols}

All R4RS predicates, accessors, and mutators are provided:
{\tt boolean?}, {\tt eq?}, {\tt eqv?}, {\tt equal?},
{\tt string?}, {\tt string-length}, {\tt string-ref}, {\tt string-set!},
{\tt substring}, {\tt string-append}, {\tt string-copy}, {\tt string-fill!},
{\tt string->list}, {\tt list->string},
{\tt string->symbol}, {\tt symbol->string},
{\tt char?}, {\tt char->integer}, {\tt integer->char}, {\tt char-ready?},
{\tt vector?}, {\tt make-vector}, {\tt vector-ref}, {\tt vector-set!},
{\tt vector-fill!}, {\tt vector->list}, {\tt list->vector},
plus all string and character comparison predicates including
case-insensitive variants.

\subsection{I/O}

\begin{lstlisting}[language=Scheme]
read  read-char  peek-char  write  display  newline
write-char  load
input-port?  output-port?  eof-object?
current-input-port  current-output-port
open-input-file  close-input-port
open-output-file  close-output-port
call-with-input-file  call-with-output-file
with-input-from-file  with-output-to-file
display-noflush  write-noflush  ; no flush variants
\end{lstlisting}

\subsection{Control}

\begin{lstlisting}[language=Scheme]
call-with-current-continuation  ; alias: call/cc
force  eval
\end{lstlisting}

{\tt eval} can take an optional environment argument.

\subsection{MScheme Extensions}\label{sec:extensions}

MScheme extends R4RS with the following procedures, grouped by category.

\subsubsection{Error Handling and Diagnostics}

\nopagebreak[4]
\begin{description}
\item[{\tt (error msg \ldots)}]
  Signal an error.  All arguments are {\tt display}ed into the
  error message.  Raises {\tt Scheme.E}.
\item[{\tt (exit)} / {\tt (exit {\it code})}]
  Exit the process immediately via {\tt Process.Exit}.
  {\it code\/} is an integer Unix exit code; it defaults to~0.
  Convention: zero indicates success, non-zero indicates failure.
\item[{\tt (set-warnings-are-errors!~b)}]
  When {\it b\/} is {\tt \#t}, warnings become errors.
\item[{\tt (enable-tracebacks!)} / {\tt (disable-tracebacks!)}]
  Enable or suppress stack tracebacks on error.
  Tracebacks are enabled by default.
\item[{\tt (set-rt-error-mapping!~b)}]
  When {\tt \#t} (default), Modula-3 runtime errors (array bounds,
  narrow failure, etc.) are caught and re-raised as {\tt Scheme.E}
  exceptions.  When {\tt \#f}, they crash the process.
\end{description}

\subsubsection{Memoization}

\nopagebreak[4]
\begin{description}
\item[{\tt (eq?-memo proc)}]
  Return a memoized version of the single-argument procedure
  {\it proc}.  The cache is keyed by {\tt eq?} identity on the
  argument, so it is suitable for memoizing over Modula-3 object
  references or symbols.
\item[{\tt (equal?-memo proc)}]
  Like {\tt eq?-memo}, but the cache is keyed by {\tt equal?}, so
  it works for strings, lists, and other structural values.
\end{description}

\noindent Example:
\begin{lstlisting}[language=Scheme]
(define slow-fib
  (lambda (n)
    (if (< n 2) n
        (+ (fast-fib (- n 1)) (fast-fib (- n 2))))))
(define fast-fib (equal?-memo slow-fib))
(fast-fib 100)  ; instant
\end{lstlisting}

\subsubsection{Formatting and Introspection}

\nopagebreak[4]
\begin{description}
\item[{\tt (fmtreal val style prec)}]
  Format a floating-point number.  {\it style\/} is one of the symbols
  {\tt 'auto}, {\tt 'fix}, or {\tt 'sci}.  {\it prec\/} is the number of
  digits after the decimal point.  Returns a Modula-3 {\tt TEXT}.
\item[{\tt (stringify obj)}]
  Convert any Scheme object to its printed representation as a
  string (like {\tt write}, but returns a string instead of writing
  to a port).
\item[{\tt (refrecord-format obj)}]
  Introspect a Modula-3 {\tt REF RECORD} value and return its
  field structure as a Scheme list.
\item[{\tt (macro-expand expr)}]
  Expand one level of macro application in {\it expr}.
\item[{\tt (string-havesub? str sub)}]
  Return {\tt \#t} if the string {\it sub\/} occurs as a substring
  of {\it str}.
\end{description}

\subsubsection{Numeric Type Conversions}

MScheme's native number type is {\tt LONGREAL}.
These procedures perform explicit narrowing for Modula-3 interop.

\nopagebreak[4]
\begin{description}
\item[{\tt (number->LONGREAL x)}]
  Identity on the numeric value (MScheme numbers are already
  {\tt LONGREAL}).  Useful as an explicit type annotation when passing
  to stubs that distinguish {\tt REAL} vs.\ {\tt LONGREAL}.
\item[{\tt (number->REAL x)}]
  Narrow {\it x\/} from {\tt LONGREAL} to {\tt REAL} (single
  precision).  May lose precision.
\item[{\tt (number->EXTENDED x)}]
  Narrow {\it x\/} to {\tt EXTENDED} (extended precision on
  platforms that support it).
\end{description}

\subsubsection{Environment Manipulation}

\nopagebreak[4]
\begin{description}
\item[{\tt (dump-environment wr)}]
  Serialize (pickle) the entire global environment to the Modula-3
  writer {\it wr\/} (a {\tt Wr.T}, not a filename).  Captures all
  top-level bindings, including closures.
\item[{\tt (load-environment! rd)}]
  Deserialize (unpickle) a global environment from the Modula-3
  reader {\it rd\/} (a {\tt Rd.T}), replacing the current global
  environment.
\item[{\tt (change-global-environment! env)}]
  Replace the interpreter's global environment with {\it env}, which
  must be an environment object (as returned by
  {\tt current-environment}).  All subsequent top-level evaluations
  use the new environment.
\item[{\tt (get-parent-environment env)}]
  Return the parent (enclosing) environment of {\it env}.  Returns
  {\tt \#f} if {\it env\/} is the top-level environment.
\item[{\tt (global-exists? sym)}]
  Return {\tt \#t} if the symbol {\it sym\/} is bound in the global
  environment.
\item[{\tt (define-global-symbol sym val)}]
  Bind {\it sym\/} to {\it val\/} in the global environment,
  regardless of the current lexical scope.
\end{description}

\subsubsection{Random Numbers and Timing}

\nopagebreak[4]
\begin{description}
\item[{\tt (random)}]
  Return a uniformly distributed random {\tt LONGREAL} in
  $[0, 1)$.
\item[{\tt (normal)} / {\tt (normal mean stddev)}]
  Return a normally distributed random number.  With no arguments,
  mean is~0 and standard deviation is~1.
\item[{\tt (time-call thunk)} / {\tt (time-call thunk label)}]
  Call {\it thunk\/} (a zero-argument procedure) and return
  its result, printing the elapsed time.  If {\it label\/} is
  provided, it prefixes the timing output.
\end{description}

\subsubsection{System Interaction}

\nopagebreak[4]
\begin{description}
\item[{\tt (system cmd)}]
  Execute the shell command {\it cmd\/} (a string) and return its
  output as a {\tt TEXT}.
\item[{\tt (read-big-int)} / {\tt (read-big-int port)}]
  Read an arbitrary-precision integer from the current input port
  or a specified port.  Returns a {\tt BigInt.T}.
\item[{\tt (jailbreak proc)}]
  Escape hatch for sandboxed interpreters.  Allows a single
  procedure {\it proc\/} to be used to add extensions to a
  locked-down {\tt Scheme.T} that was created with
  {\tt SandboxDefiner}.
\end{description}

\subsubsection{Low-Level Modula-3 Interop}

\nopagebreak[4]
\begin{description}
\item[{\tt (modula-3-op name obj)}]
  Low-level method dispatch: invoke the method {\it name\/}
  (a symbol) on the Modula-3 object {\it obj}.
\item[{\tt (encode-pickle-vector obj)}]
  Serialize a Modula-3 object to a byte vector ({\tt REF ARRAY OF
  CHAR}) using the M3 pickle protocol.
\item[{\tt (decode-pickle-vector vec)}]
  Deserialize a Modula-3 object from a byte vector previously
  produced by {\tt encode-pickle-vector}.
\end{description}

\subsubsection{Debugging}

\nopagebreak[4]
\begin{description}
\item[{\tt (debug-setenv name value)}]
  Set the Debug module environment variable {\it name\/} to
  {\it value}.
\item[{\tt (debug-clearenv name)}]
  Clear the Debug module environment variable {\it name}.
\item[{\tt (debug-haveenv name)}]
  Return {\tt \#t} if the Debug environment variable {\it name\/}
  is set.
\item[{\tt (stdio-stderr)}]
  Return the process's standard error stream as a {\tt Wr.T}.
  Useful for directing debug output separately from
  {\tt current-output-port}.
\end{description}


%======================================================================
\chapter{Bundled Scheme Libraries}\label{ch:libraries}
%======================================================================

Libraries are bundled into the MScheme executable via {\tt m3bundle}
and can be loaded by name without a file extension.  They reside in\\
{\tt m3-scheme/scheme-lib/src/}.

\section{{\tt require} --- Module System}

Loaded automatically by {\tt SchemeM3.T}.  Provides {\tt require-modules}
(see Section~\ref{sec:modules}).

\section{{\tt basic-defs} --- Utility Functions}

\begin{lstlisting}[language=Scheme]
(require-modules "basic-defs")
\end{lstlisting}

\noindent Provides:
\begin{itemize}
\item {\tt (nth lst n)} --- zero-indexed list access
\item {\tt (mem? eq-fn x lst)} --- generic membership test
\item {\tt (member? x lst)} --- {\tt equal?}-based membership
\item {\tt (filter pred lst)} --- filter list by predicate
\item {\tt (accumulate op init seq)} --- fold right (SICP-style)
\item {\tt (identity x)} --- identity function
\item {\tt (last lst)}, {\tt (all-except-last lst)}
\item {\tt (head n lst)}, {\tt (tail n lst)} --- first/last $n$ elements
\item {\tt (uniq eq-fn lst)} --- remove duplicates
\item {\tt (map2 op-normal op-last lst)} --- map with different op for last
\end{itemize}

Also loads {\tt basic-mbe} (hygienic macro support primitives).

\section{{\tt display} --- Output Helpers}

\begin{lstlisting}[language=Scheme]
(require-modules "display")
\end{lstlisting}

\noindent Provides:
\begin{itemize}
\item {\tt dnl} --- the newline character ({\tt \#\textbackslash newline}),
  used as a separator
\item {\tt (dis arg1 arg2 \ldots\ [port])} --- display multiple items,
  optionally to a given port
\end{itemize}

\noindent Example:
\begin{lstlisting}[language=Scheme]
(dis "Hello, " name "!" dnl)
(dis "x = " x ", y = " y dnl)
\end{lstlisting}

\section{{\tt hashtable} --- Hash Tables}

\begin{lstlisting}[language=Scheme]
(require-modules "hashtable")
\end{lstlisting}

Creates message-passing hash table objects:

\begin{lstlisting}[language=Scheme]
(define tbl (make-hash-table 100 my-hash-fn))
(tbl 'add-entry! "key" "value")
(tbl 'retrieve "key")          ; => "value"
(tbl 'update-entry! "key" "new-value")
(tbl 'delete-entry! "key")
(tbl 'keys)                    ; => list of all keys
(tbl 'values)                  ; => list of all values
(tbl 'size)                    ; => number of entries
(tbl 'clear!)                  ; remove all entries
(tbl 'rehash! new-size)        ; resize
\end{lstlisting}

\noindent Convenience constructor:
\begin{lstlisting}[language=Scheme]
(define tbl (make-string-hash-table 100))
\end{lstlisting}

\noindent The sentinel {\tt '*hash-table-search-failed*} is returned
for missing keys.  Test with {\tt eq?}:
\begin{lstlisting}[language=Scheme]
(if (not (eq? (tbl 'retrieve k)
              '*hash-table-search-failed*))
    (tbl 'retrieve k)
    default-value)
\end{lstlisting}

\section{{\tt set} --- Sets}

\begin{lstlisting}[language=Scheme]
(require-modules "set")
\end{lstlisting}

Creates message-passing set objects built on hash tables:

\begin{lstlisting}[language=Scheme]
(define s (make-symbol-set 100))
(s 'insert! 'apple)       ; returns #f if new
(s 'member? 'apple)       ; => #t
(s 'delete! 'apple)
(s 'keys)                 ; list of elements
(s 'size)
(s 'copy)
(s 'union other-set)
(s 'intersection other-set)
(s 'diff other-set)       ; set difference
\end{lstlisting}

Also provides {\tt make-symbol-hash-table} and {\tt make-string-set}.

\section{{\tt struct} --- Record Types}

\begin{lstlisting}[language=Scheme]
(require-modules "struct")
\end{lstlisting}

Defines record-like structured types with message-passing access:

\begin{lstlisting}[language=Scheme]
(define point-type
  (make-struct-type 'point
    `((x 0) (y 0) (label "origin"))))

(define p (point-type 'new))
(p 'get 'x)              ; => 0
(p 'set! 'x 42)
(p 'get 'x)              ; => 42
(p 'inc! 'x)             ; increment by 1
(p 'inc! 'x 10)          ; increment by 10
(p 'inc! 'x (lambda (v) (* v 2)))
(p 'display)             ; => ((x . 106) (y . 0) ...)
(p 'copy)                ; deep copy
(p 'set 'x 99)           ; functional update (new copy)
(p 'type)                ; => the type object
(p 'reset!)              ; reset to initial values
\end{lstlisting}

Dynamic initializers: if a field's default is a procedure, it is called
(with no arguments) to produce each new value:

\begin{lstlisting}[language=Scheme]
(make-struct-type 'thing
  `((id  ,(lambda () (random)))
    (data ())))
\end{lstlisting}

Retrieve a type by name with {\tt (get-struct-type 'point)}.

\section{{\tt m3} --- Modula-3 Bindings}\label{sec:m3lib}

\begin{lstlisting}[language=Scheme]
(require-modules "m3")
\end{lstlisting}

This is the critical library for Modula-3 interoperation.  It installs
all registered Scheme procedure stubs as top-level Scheme bindings
using the convention {\tt Interface.Procedure}:

\begin{lstlisting}[language=Scheme]
; Before "m3": must use raw stub call
(scheme-procedure-stubs-call '(Fmt . Int) (list 42) '())

; After "m3": direct call
(Fmt.Int 42)
\end{lstlisting}

Also provides the key helper:
\begin{lstlisting}[language=Scheme]
(obj-method-wrap obj 'TypeName.T)
\end{lstlisting}

And introspection utilities {\tt schemify}, {\tt have-type-ops?},
{\tt closest-opped-supertype}, {\tt closest-type-op}.

\section{Other Libraries}

\begin{itemize}
\item {\tt mergesort} --- merge sort for lists
\item {\tt mbe} --- R5RS hygienic macros ({\tt syntax-rules})
\item {\tt time} --- time formatting utilities
\item {\tt exit} --- exit handlers
\item {\tt gnuplot} --- gnuplot integration
\item {\tt history} --- command history
\item {\tt scodegen} --- SQL schema code generation (Hibernate-inspired)
\end{itemize}


%======================================================================
\chapter{Modula-3 Integration}\label{ch:m3interop}
%======================================================================

This is the heart of MScheme's value.  The system provides transparent
bidirectional bindings between Scheme and Modula-3.

\section{Zero-Copy Value Sharing}

Because {\tt SchemeObject.T = REFANY}, and all Modula-3 heap objects
are {\tt REFANY}, \emph{Modula-3 objects are Scheme values without any
conversion.}  When a stub-generated procedure returns a Modula-3
object, the exact same heap pointer is the Scheme value.  When Scheme
passes that value back to Modula-3, no conversion occurs.

Only Modula-3 value types ({\tt INTEGER}, {\tt REAL}, {\tt BOOLEAN},
{\tt CHAR}, records, arrays) require boxing and unboxing at the
boundary, and this is handled automatically by the generated stubs.

\section{The Stub Generator ({\tt sstubgen})}\label{sec:sstubgen}

\subsection{How It Works}

{\tt sstubgen} is a build-time tool that reads Modula-3 {\tt .i3}
interface files via the M3 Toolkit ({\tt m3tk})~\cite{m3tk} and generates:

\begin{enumerate}
\item \textbf{Type converters} --- {\tt ToModula} (Scheme$\to$M3, with
  type and bounds checking) and {\tt ToScheme} (M3$\to$Scheme) for
  each type.
\item \textbf{Procedure call stubs} --- wrappers that unpack Scheme
  argument lists, convert types, call the M3 procedure, and convert
  the result.
\item \textbf{Object method dispatch} --- for M3 object types,
  generates method dispatch tables accessible via {\tt modula-type-op}.
\item \textbf{Constructor registration} --- {\tt new-modula-object}
  support for object types.
\item \textbf{Type metadata} --- runtime type information (fields,
  methods, supertypes) pickled and registered for introspection.
\end{enumerate}

The {\tt sstubgen} tool shares a design lineage with the DEC~SRC
network objects stub generator ({\tt stubgen})~\cite{netobj}, both
based on the M3 Toolkit.

\subsection{Build System Integration}

Stub generation is automated by the Quake template
{\tt schemestubs.tmpl}, which defines the following build directives:

\begin{center}
\begin{tabular}{ll}
\toprule
\textbf{Directive} & \textbf{Description} \\
\midrule
{\tt SchemeStubs("Fmt")}           & Generate stubs for {\tt Fmt.i3} \\
{\tt Smodule("Fmt")}               & {\tt Module} + {\tt SchemeStubs} \\
{\tt Sinterface("Fmt")}            & {\tt Interface} + {\tt SchemeStubs} \\
{\tt Ssequence("K","V")}           & {\tt Sequence} + {\tt SchemeStubs("KSeq")} \\
{\tt ExportSchemeStubs("prog")}    & Export stubs for linking (phase~2) \\
{\tt importSchemeStubs()}          & Import stubs from dependencies (phase~3) \\
\bottomrule
\end{tabular}
\end{center}

\subsection{Three-Phase Stub Generation}

\begin{enumerate}
\item \textbf{Phase~1} ({\tt make-standard-stuff}): For each {\tt
  SchemeStubs("Fmt")} call, {\tt sstubgen} processes {\tt Fmt.i3} and
  generates {\tt FmtSchemeStubs.i3} and {\tt FmtSchemeStubs.m3}.
  These register procedures, types, and operations at module
  initialization time.

\item \textbf{Phase~2} ({\tt write-scheme-package-exports}): For each
  {\tt ExportSchemeStubs("prog")}, generates an interface that lists
  all stub modules in this package.

\item \textbf{Phase~3} ({\tt importSchemeStubs}): Pulls in stub
  exports from all imported packages, ensuring all stubs are linked
  into the executable.
\end{enumerate}

\subsection{Worked Example: Stub Generation for {\tt Fmt}}\label{sec:fmtexample}

This section traces the complete stub-generation pipeline for
{\tt Fmt.i3}, the standard Modula-3 formatting interface.

\paragraph{Step 1: The Modula-3 interface.}
The relevant declarations in {\tt Fmt.i3}:

\begin{lstlisting}[language=Modula3]
INTERFACE Fmt;

TYPE Base = [2..16];
TYPE Style = {Sci, Fix, Auto};
TYPE Align = {Left, Right};

PROCEDURE Bool(b: BOOLEAN): TEXT;
PROCEDURE Int(n: INTEGER; base: Base := 10): TEXT;
PROCEDURE LongReal(x: LONGREAL; style := Style.Auto;
    prec: CARDINAL := 15; literal := FALSE): TEXT;
PROCEDURE Pad(text: TEXT; length: CARDINAL;
    padChar: CHAR := ' '; align: Align := Align.Right): TEXT;
\end{lstlisting}

\paragraph{Step 2: The {\tt m3makefile} directive.}
A single line causes stubs to be generated at build time:

\begin{lstlisting}[language={}]
SchemeStubs ("Fmt")
\end{lstlisting}

\paragraph{Step 3: The generated code.}
{\tt sstubgen} produces {\tt FmtSchemeStubs.m3} containing three
kinds of generated code:

\medskip\noindent\textbf{(a) Procedure call stubs.}\  Each M3 procedure gets a wrapper
that unpacks a Scheme argument list, converts types, calls the M3
procedure, and converts the result:

\begin{lstlisting}[language=Modula3]
PROCEDURE CallStub_Fmt_Int(
    interp : Scheme.T;
    args : SchemeObject.T;
    excHandler : SchemeObject.T) : SchemeObject.T
  RAISES { Scheme.E } =
  PROCEDURE Next() : SchemeObject.T RAISES { Scheme.E } =
    BEGIN
      TRY RETURN SchemeUtils.First(p__)
      FINALLY p__ := SchemePair.Pair(
                       SchemeUtils.Rest(p__))
      END
    END Next;
  VAR p__ := SchemePair.Pair(args);
  BEGIN
    TRY
      VAR formal_n :=
            SchemeModula3Types.ToModula_INTEGER(Next());
          formal_base := ToModula_Fmt_Base(Next());
      BEGIN
        WITH res = SchemeModula3Types.ToScheme_TEXT(
                     Fmt.Int(formal_n, formal_base)) DO
          RETURN res
        END
      END
    EXCEPT END;
    RETURN SchemeBoolean.False()
  END CallStub_Fmt_Int;
\end{lstlisting}

\noindent The key pattern: {\tt Next()} walks the Scheme argument
list, {\tt ToModula\_*} functions convert each argument, and
{\tt ToScheme\_TEXT} wraps the result.

\medskip\noindent\textbf{(b) Type converters.}\  Subranges are validated at the
boundary:

\begin{lstlisting}[language=Modula3]
PROCEDURE ToModula_Fmt_Base(x : SchemeObject.T)
    : Fmt.Base  RAISES { Scheme.E } =
  BEGIN
    WITH baseVal =
           SchemeModula3Types.ToModula_INTEGER(x) DO
      IF baseVal < FIRST(Fmt.Base)
      OR baseVal > LAST(Fmt.Base) THEN
        RAISE Scheme.E("Value out of range for "
          & "Fmt.Base :" & SchemeUtils.Stringify(x))
      END;
      RETURN baseVal
    END
  END ToModula_Fmt_Base;
\end{lstlisting}

\noindent Enumerations are mapped to Scheme symbols:

\begin{lstlisting}[language=Modula3]
PROCEDURE ToModula_Fmt_Style(x : SchemeObject.T)
    : Fmt.Style  RAISES { Scheme.E } =
  BEGIN
    CONST
      Name = ARRAY Fmt.Style OF TEXT
               { "Sci", "Fix", "Auto" };
    BEGIN
      FOR i := FIRST(Name) TO LAST(Name) DO
        IF SchemeSymbol.FromText(Name[i]) = x THEN
          RETURN i
        END
      END;
      RAISE Scheme.E("Not a value of Fmt.Style : "
        & SchemeUtils.Stringify(x))
    END
  END ToModula_Fmt_Style;
\end{lstlisting}

\medskip\noindent\textbf{(c) Registration.}\  A {\tt RegisterStubs}
procedure, called at module initialization, registers every stub
with the Scheme runtime:

\begin{lstlisting}[language=Modula3]
PROCEDURE RegisterStubs() =
  BEGIN
    SchemeProcedureStubs.Register(
      NEW(SchemeProcedureStubs.Qid,
        intf := Atom.FromText("Fmt"),
        item := Atom.FromText("Int")),
      CallStub_Fmt_Int);
    (* ... one call per procedure ... *)
  END RegisterStubs;
\end{lstlisting}

\paragraph{Step 4: Using from Scheme.}
After {\tt (require-modules "m3")}, all registered stubs are
available:

\begin{lstlisting}[language=Scheme]
(require-modules "m3")
(Fmt.Int 42)                    ; => "42"
(Fmt.Int 255 16)                ; => "ff"
(Fmt.Bool #t)                   ; => "TRUE"
(Fmt.LongReal 3.14159 'Fix 2 #f)  ; => "3.14"
(Fmt.Pad "hi" 10 #\space 'Left)   ; => "hi        "
\end{lstlisting}

\paragraph{Key observations.}
\begin{itemize}
\item Enumerations ({\tt Style}, {\tt Align}) become Scheme symbols:
  {\tt 'Sci}, {\tt 'Fix}, {\tt 'Auto}, {\tt 'Left}, {\tt 'Right}.
\item Subranges ({\tt Base = [2..16]}) are validated at the stub
  boundary; passing~1 or~17 raises a Scheme error.
\item {\tt TEXT} $\leftrightarrow$ Scheme string conversion is
  transparent---no explicit marshaling needed.
\item Default arguments are flattened by {\tt sstubgen}: all
  parameters must be provided explicitly from Scheme.
  For example, {\tt Fmt.LongReal} requires four arguments, not one.
\end{itemize}


\subsection{Example {\tt m3makefile}}

\begin{lstlisting}[language={}]
import ("mscheme")
import ("modula3scheme")
import ("sstubgen")

Smodule ("MyCompiler")         % Module + SchemeStubs
SchemeStubs ("MyAST")          % stubs only
SchemeStubs ("Text")           % standard library
SchemeStubs ("Fmt")

ExportSchemeStubs ("myapp")
importSchemeStubs()

implementation ("Main")
program ("myapp")
\end{lstlisting}

\section{Using Modula-3 Bindings from Scheme}\label{sec:usingm3}

After loading the {\tt "m3"} library
(Section~\ref{sec:m3lib}), all stub-registered procedures become
available as {\tt InterfaceName.ProcedureName}:

\begin{lstlisting}[language=Scheme]
(require-modules "m3")

;; Call Fmt.Int from M3's standard library
(Fmt.Int 42)              ; => "42"
(Fmt.Int 255 16)          ; => "ff"

;; String operations
(Text.Length "hello")     ; => 5
(Text.Sub "hello" 1 3)   ; => "ell"

;; Environment variables
(Env.Get "HOME")          ; => "/home/user"

;; File I/O through M3
(define wr (FileWr.Open "/tmp/test.txt"))
(Wr.PutText wr "Hello, world!\n")
(Wr.Close wr)
\end{lstlisting}

\section{Working with Modula-3 Objects}

\subsection{Creating Objects}

\begin{lstlisting}[language=Scheme]
(define obj (new-modula-object 'MyClass.T))
\end{lstlisting}

The type name uses the {\tt Interface.Type} convention.
Initialization arguments may be passed as additional arguments.

\subsection{Method Calls via {\tt obj-method-wrap}}

The {\tt obj-method-wrap} function creates a Scheme closure that
dispatches method calls on a Modula-3 object:

\begin{lstlisting}[language=Scheme]
(define wrapped (obj-method-wrap raw-obj 'TypeName.T))
(wrapped 'methodName arg1 arg2 ...)
\end{lstlisting}

The special symbol {\tt '*m3*} retrieves the underlying Modula-3 object:
\begin{lstlisting}[language=Scheme]
(wrapped '*m3*)  ; => the raw M3 REFANY
\end{lstlisting}

\subsection{Complete Example}

\begin{lstlisting}[language=Scheme]
;; Create a text sequence (TextSeq.T from libm3)
(define seq (obj-method-wrap
              (new-modula-object 'TextSeq.T)
              'TextSeq.T))
(seq 'addhi "hello")
(seq 'addhi "world")
(seq 'size)           ; => 2
(seq 'get 0)          ; => "hello"
(seq 'get 1)          ; => "world"
\end{lstlisting}

\subsection{Convenience Pattern: {\tt init-seq}}

A common pattern in applications wraps sequence creation:

\begin{lstlisting}[language=Scheme]
(define (init-seq type-name)
  (let ((res (obj-method-wrap
               (new-modula-object type-name)
               type-name)))
    (obj-method-wrap (res 'init 10) type-name)))

(define stmts (init-seq 'CspStatementSeq.T))
(stmts 'addhi some-statement)
(stmts '*m3*)     ; get the raw M3 sequence
\end{lstlisting}

\section{Low-Level Stub API}

These primitives underlie the {\tt "m3"} library and are available
even without loading it:

\begin{lstlisting}[language=Scheme]
;; List all registered procedure stubs
(scheme-procedure-stubs-list)
; => ((Fmt . Int) (Fmt . LongReal) (Text . Length) ...)

;; Call a stub directly
(scheme-procedure-stubs-call '(Fmt . Int) (list 42) '())

;; List all types that can be "new"ed
(newable-types)

;; Create a new M3 object
(new-modula-object 'TypeName.T)

;; Call an operation on an M3 object by typecode
(modula-type-op tc-or-symbol 'op-name object args ...)

;; List operations for a typecode
(list-modula-type-ops typecode)

;; Type system introspection
(list-modula-types)
(typename->typecode 'TypeName.T)
(rttype-typecode obj)
(rttype-maxtypecode)
(rttype-supertype tc)
(rttype-issubtype? tc1 tc2)
(rttype-ndimensions tc)
(rtbrand-getname tc)
(modula-type-class 'TypeName.T)
\end{lstlisting}

\section{{\tt SchemeM3.T} Extended Primitives}

When using {\tt SchemeM3.T} (the standard extended interpreter), these
additional primitives are available:

\begin{lstlisting}[language=Scheme]
;; System utilities
(gc)                          ; force GC
(timenow)                     ; Unix epoch time
(time->string time [tz])      ; "2024-MAR-14 10:30:00.000"
(time->list time [tz])        ; (year month day wday hr min sec)
(hostname)                    ; machine hostname (returns symbol)
(getunixpid)                  ; Unix PID

;; Formatting
(stringify obj)               ; like write, returns string
(fmtreal value 'style prec)  ; 'auto, 'fix, or 'sci

;; File operations
(filewr-open filename)
(filewr-open 'append filename)
(wr-close wr)
(remove-file path)
(fs-rename old new)
(cmp-files f1 f2)             ; binary compare
(fs-status-modificationtime path)

;; Debugging
(debug-setlevel n)
(debug-addstream wr)
(debug-remstream wr)
(stdio-stderr)

;; Network Objects
(netobj-export-global-environment "name")
\end{lstlisting}

\section{Asynchronous Execution}

MScheme provides thread-based asynchronous execution:

\begin{lstlisting}[language=Scheme]
;; Schedule procedure to run at a future time
(define (my-func) (+ 40 2))
(define h (at-run (+ (timenow) 5.0) my-func callback))
(at-join h)               ; wait for result => 42

;; Periodic execution (every 10 seconds)
(define (check-something) (display ".") #t)
(define clk (clock-run 10.0 check-something err-handler))
;; runs until procedure returns #f
(clock-stop clk)
(clock-list)               ; list active clocks
\end{lstlisting}

\section{Running Shell Commands}

\begin{lstlisting}[language=Scheme]
(run-command "ls" "-la" "/tmp")
(run-command-with-timeout 30 "make" "-j4")
(system "echo hello")     ; returns output as text
\end{lstlisting}


%======================================================================
\chapter{BigInt Support}
%======================================================================

MScheme integrates with the {\tt BigInt} module for
arbitrary-precision integer arithmetic.  After loading stubs:

\begin{lstlisting}[language=Scheme]
(define a (BigInt.New 12345678901234567890))
(define b (BigInt.New 98765432109876543210))

(BigInt.Add a b)
(BigInt.Mul a b)
(BigInt.Div a b)
(BigInt.Mod a b)
(BigInt.Compare a b)      ; => -1, 0, or 1
(BigInt.Equal a b)        ; => #t or #f
(BigInt.Format a 10)      ; => "12345678901234567890"
(BigInt.Format a 16)      ; hex string
(BigInt.ToInteger a)      ; => Scheme number (if fits)
(BigInt.IsT x)            ; => #t if BigInt
(BigInt.Shift a n)        ; left shift
(BigInt.And a b)           (BigInt.Or a b)
(BigInt.Xor a b)           (BigInt.Not a)
(BigInt.GetAbsMsb a)      ; highest set bit position
(BigInt.ScanBased "ff" 16 #f)  ; parse string
\end{lstlisting}

BigInts can be read from input ports:
\begin{lstlisting}[language=Scheme]
(read-big-int)            ; from current input
(read-big-int port)       ; from specific port
\end{lstlisting}


%======================================================================
\chapter{Wrapping C Libraries}\label{ch:clibrary}
%======================================================================

One of MScheme's most powerful capabilities is making C libraries
accessible from Scheme with full garbage-collected memory management
and no unsafe code visible to the Scheme programmer.  This chapter
uses the GMP (GNU Multiple Precision) integer library as a
complete, working example.

\section{The Pattern}

The integration follows a layered architecture:

\begin{center}
\begin{tabular}{ccc}
\textbf{Layer} & & \textbf{Artifact} \\
\midrule
C library & $\longrightarrow$ & {\tt libgmp.a} \\
C helpers & $\longrightarrow$ & {\tt MpzC.c} \\
{\tt <*EXTERNAL*>} bindings & $\longrightarrow$ & {\tt MpzP.i3} \\
Safe M3 wrapper + WeakRef & $\longrightarrow$ & {\tt Mpz.m3} \\
Type revelation & $\longrightarrow$ & {\tt MpzRep.i3} \\
{\tt SchemeStubs("Mpz")} & $\longrightarrow$ & {\tt MpzSchemeStubs.m3} \\
Scheme code & $\longrightarrow$ & {\tt (Mpz.add c a b)} \\
\end{tabular}
\end{center}

\noindent Each layer adds safety.  The Scheme programmer at the top
sees only safe, garbage-collected objects.

\section{Layer 1: C Interface}

{\tt MpzP.i3} declares Modula-3 bindings to C functions using the
{\tt <*EXTERNAL*>} pragma.  Each declaration names the C symbol
explicitly:

\begin{lstlisting}[language=Modula3]
INTERFACE MpzP;
IMPORT Word, Ctypes;

TYPE MpzPtrT = ADDRESS;

<*EXTERNAL "__gmpz_init"*>
PROCEDURE c_init(f0: MpzPtrT);

<*EXTERNAL "__gmpz_clear"*>
PROCEDURE c_clear(f0: MpzPtrT);

<*EXTERNAL "__gmpz_add"*>
PROCEDURE c_add(f0: MpzPtrT; f1: MpzPtrT;
                f2: MpzPtrT);

<*EXTERNAL "__gmpz_set_si"*>
PROCEDURE c_set_si(f0: MpzPtrT; f1: INTEGER);

<*EXTERNAL "__gmpz_cmp"*>
PROCEDURE c_cmp(f0: MpzPtrT; f1: MpzPtrT)
    : Ctypes.int;

(* formatting helpers -- implemented in MpzC.c *)
<*EXTERNAL mpz_format_decimal*>
PROCEDURE format_decimal(z: MpzPtrT)
    : Ctypes.const_char_star;

<*EXTERNAL mpz_free_formatted*>
PROCEDURE free_formatted(f0: Ctypes.char_star);
\end{lstlisting}

\noindent The C helper {\tt MpzC.c} provides formatting functions
that GMP does not directly expose:

\begin{lstlisting}[language={}]
/* MpzC.c */
#include <gmp.h>

char *mpz_format_decimal(const mpz_t z) {
  char *p;
  gmp_asprintf(&p, "%Zd", z);
  return p;
}

void mpz_free_formatted(char *p) {
  void (*freefunc)(void *, size_t);
  mp_get_memory_functions(NULL, NULL, &freefunc);
  freefunc(p, 0);
}
\end{lstlisting}

\section{Layer 2: Safe Wrapper with Garbage Collection}

The key module is {\tt Mpz.m3}.  It wraps every C operation in a
safe Modula-3 procedure and---crucially---uses {\tt WeakRef} to
ensure that GMP memory is freed when the Modula-3 garbage collector
reclaims the object:

\begin{lstlisting}[language=Modula3]
UNSAFE MODULE Mpz;
IMPORT MpzRep, MpzP AS P, WeakRef, M3toC;

PROCEDURE New() : T =
  VAR res := NEW(T);
  BEGIN
    P.c_init(LOOPHOLE(ADR(res.val), P.MpzPtrT));
    EVAL WeakRef.FromRef(res, CleanUp);
    set_ui(res, 0);
    RETURN res
  END New;

PROCEDURE CleanUp(
    <*UNUSED*> READONLY w: WeakRef.T;
    r: REFANY) =
  BEGIN
    WITH this = NARROW(r, T) DO
      P.c_clear(ADR(this.val))
    END
  END CleanUp;

PROCEDURE FormatDecimal(t: T) : TEXT =
  VAR cs := P.format_decimal(ADR(t.val));
  BEGIN
    TRY
      RETURN M3toC.CopyStoT(cs)
    FINALLY
      P.free_formatted(cs)
    END
  END FormatDecimal;
\end{lstlisting}

\noindent The critical line is:
\begin{lstlisting}[language=Modula3]
EVAL WeakRef.FromRef(res, CleanUp);
\end{lstlisting}

\noindent This registers {\tt CleanUp} as a weak-reference callback.
When the garbage collector determines that the {\tt Mpz.T} object is
unreachable, it schedules a call to {\tt CleanUp}, which calls
{\tt gmpz\_clear} to free the C-allocated memory.  The Scheme (or
Modula-3) programmer never needs to call a ``free'' or ``dispose''
function.

\paragraph{{\tt WeakRef} API summary.}

\begin{lstlisting}[language=Modula3]
(* WeakRef.i3 *)
PROCEDURE FromRef(r: REFANY;
    p: CleanUpProc := NIL): T;
PROCEDURE ToRef(w: T): REFANY;
TYPE CleanUpProc =
    PROCEDURE(READONLY w: T; r: REFANY);
\end{lstlisting}

\noindent {\tt FromRef} creates a weak reference and associates
a cleanup procedure.  {\tt ToRef} dereferences a weak reference,
returning {\tt NIL} if the referent has been collected.

\section{Layer 3: Type Revelation}

{\tt MpzRep.i3} reveals the opaque type {\tt Mpz.T} as a branded
object containing the C {\tt mpz\_t} struct inline:

\begin{lstlisting}[language=Modula3]
INTERFACE MpzRep;
IMPORT Mpz, Word;

REVEAL
  Mpz.T = BRANDED Mpz.Brand OBJECT
    val : Val;
  END;

TYPE
  Val = RECORD
    w : ARRAY [0..2] OF Word.T;
    (* cast to mpz_t -- at least big enough *)
  END;
END MpzRep.
\end{lstlisting}

\noindent The three-word array is large enough to hold the GMP
{\tt mpz\_t} structure on all supported platforms.  The {\tt BRANDED}
keyword ensures type safety---different {\tt Mpz.T} values cannot be
confused with other {\tt REFANY} values at runtime.

\section{Layer 4: Build Integration}

The {\tt m3makefile} for the {\tt mpz} package:

\begin{lstlisting}[language={}]
import ("libm3")
import ("m3core")

if SYSTEM_LIBS contains "MPZ"
  import_sys_lib ("MPZ")
  c_source ("MpzC")
end

Interface ("MpzP")
interface ("MpzRep")
Module ("Mpz")
library ("mpz")
\end{lstlisting}

\noindent The conditional {\tt if SYSTEM\_LIBS contains "MPZ"}
allows the package to build without GMP installed (the C source
is simply skipped).  The top-level {\tt m3overrides} file defines:
\begin{lstlisting}[language={}]
SYSTEM_LIBS{"MPZ"} = [ "-lgmp" ]
\end{lstlisting}

\section{Layer 5: Scheme Stub Generation}

In the application that wants to expose {\tt Mpz} to Scheme, the
{\tt m3makefile} adds:

\begin{lstlisting}[language={}]
import ("mpz")
SchemeStubs ("Mpz")
ExportSchemeStubs ("myapp")
importSchemeStubs ()
\end{lstlisting}

\noindent {\tt sstubgen} reads {\tt Mpz.i3} and generates call stubs
for every procedure: {\tt New}, {\tt add}, {\tt mul}, {\tt cmp},
{\tt FormatDecimal}, etc.

\section{Layer 6: Using from Scheme}

\begin{lstlisting}[language=Scheme]
(require-modules "m3")

;; Create three GMP integers
(define a (Mpz.New))
(define b (Mpz.New))
(define c (Mpz.New))

;; Set values
(Mpz.set_si a 42)
(Mpz.set_si b 100)

;; Arithmetic (result stored in first argument)
(Mpz.add c a b)
(Mpz.FormatDecimal c)       ; => "142"

;; More operations
(Mpz.mul c a b)
(Mpz.FormatDecimal c)       ; => "4200"

(Mpz.pow_ui c a 10)
(Mpz.FormatDecimal c)       ; => "17490122876983048192"

;; Comparison
(Mpz.cmp a b)               ; => -1  (a < b)

;; When a, b, c become unreachable, the GC calls
;; CleanUp which calls gmpz_clear -- no manual
;; memory management needed.
\end{lstlisting}

\section{Why This Matters}

\begin{center}
\fbox{\parbox{0.85\textwidth}{
\textbf{Key insight:}\  The Scheme programmer manipulates C library
objects (GMP arbitrary-precision integers) without any awareness of
C memory management.  The Modula-3 garbage collector, via
{\tt WeakRef}, automatically calls {\tt gmpz\_clear} when objects
become unreachable.  The Scheme programmer sees only safe,
garbage-collected values.

\medskip
This pattern generalizes to any C library.  The recipe is:
\begin{enumerate}
\item Declare {\tt <*EXTERNAL*>} bindings in a private {\tt .i3}
\item Write a safe wrapper module that calls {\tt WeakRef.FromRef}
  in the constructor and {\tt c\_clear}/{\tt c\_free} in the
  cleanup procedure
\item Reveal the opaque type to contain the C struct
\item Run {\tt SchemeStubs} on the safe public interface
\end{enumerate}

\noindent The result is that Scheme inherits the full power of the C
library with the safety of garbage-collected Modula-3.
}}
\end{center}


%======================================================================
\chapter{Embedding MScheme in Applications}
%======================================================================

\section{Basic Embedding}

\begin{lstlisting}[language=Modula3]
IMPORT Scheme, SchemeM3, SchemeStubs;
IMPORT SchemeReadLine, ReadLine;

BEGIN
  (* Register all auto-generated stubs *)
  SchemeStubs.RegisterStubs();

  (* Create interpreter with startup files *)
  WITH scm = NEW(SchemeM3.T).init(
    ARRAY OF TEXT { "setup.scm" }^) DO

    (* Enter interactive REPL *)
    SchemeReadLine.MainLoop(
      NEW(ReadLine.Default).init(), scm)
  END
END
\end{lstlisting}

\section{Programmatic Evaluation}

\begin{lstlisting}[language=Modula3]
(* Evaluate an expression *)
VAR result := scm.loadEvalText("(+ 1 2 3)");
(* result is SchemeObject.T = REFANY *)

(* Extract a number *)
VAR n := SchemeLongReal.FromO(result);  (* 6.0d0 *)

(* Bind a value into Scheme *)
scm.bind("my-var", SchemeLongReal.FromI(42));

(* Load and evaluate a file *)
EVAL scm.loadFile(SchemeString.FromText("mycode.scm"));
\end{lstlisting}

\section{Defining Custom Primitives}

\begin{lstlisting}[language=Modula3]
IMPORT SchemePrimitive, SchemeProcedure, SchemeM3;

(* Get the extensible primitive definer *)
VAR prims := SchemeM3.GetPrims();

(* Add a custom primitive *)
prims.addPrim("my-add",
  NEW(SchemeProcedure.T, apply := MyAddApply),
  2, 2);  (* exactly 2 args *)

PROCEDURE MyAddApply(
    <*UNUSED*> p: SchemeProcedure.T;
    <*UNUSED*> interp: Scheme.T;
    args: Scheme.Object): Scheme.Object
  RAISES { Scheme.E } =
BEGIN
  WITH x = SchemeLongReal.FromO(SchemeUtils.First(args)),
       y = SchemeLongReal.FromO(SchemeUtils.Second(args)) DO
    RETURN SchemeLongReal.FromLR(x + y)
  END
END MyAddApply;
\end{lstlisting}

\section{Security: The Sandbox}

MScheme provides three levels of primitive access:

\begin{enumerate}
\item \textbf{SandboxDefiner} --- No file I/O, no system access.  Safe
  for untrusted code.
\item \textbf{DefaultDefiner} --- Full file I/O and extensions.
\item \textbf{ExtDefiner} --- All of the above plus user-defined primitives.
\end{enumerate}

\begin{lstlisting}[language=Modula3]
(* Create a sandboxed interpreter *)
VAR scm := NEW(Scheme.T);
scm.setPrimitives(
  NEW(SchemePrimitive.SandboxDefiner));
EVAL scm.init(ARRAY OF TEXT {}, NIL);
\end{lstlisting}


%======================================================================
\chapter{The MScheme Compiler}\label{ch:compiler}
%======================================================================

MScheme includes a Level~1 ahead-of-time compiler that translates
Scheme source files into Modula-3 source code.  The compiled code
is linked into your application and executes significantly faster than
the interpreter---from 16$\times$ on tree-recursive benchmarks to over
300$\times$ on tight numeric loops, thanks to primitive inlining,
intra-module direct calls, and numeric unboxing optimizations.
The compiler preserves redefinition semantics: compiled functions look
up their callees through live environment bindings, so \texttt{set!}\
and interpreter-side \texttt{define} continue to work correctly.

The compiler is described in detail in the companion document
\emph{MScheme Compiler Design: From Binding Optimization to
Native Code Generation}.  This chapter covers practical usage.

\section{Quick Start}

Suppose you have a Scheme library \texttt{mathlib.scm}:

\begin{lstlisting}[language=Scheme]
(require-modules "basic-defs")

(define (factorial n)
  (if (= n 0) 1 (* n (factorial (- n 1)))))

(define (fibonacci n)
  (if (< n 2) n (+ (fibonacci (- n 1))
                    (fibonacci (- n 2)))))
\end{lstlisting}

To compile it, add these lines to your \texttt{m3makefile}:

\begin{lstlisting}[language={}]
import("schemecompiler")

scheme("mathlib")
\end{lstlisting}

That's it.  At build time, \texttt{cm3} invokes the MScheme compiler
to produce \texttt{mathlib.i3} and \texttt{mathlib.m3} as derived
sources.  At runtime, the compiled module self-registers in a global
table.  When Scheme code executes
\texttt{(require-modules "mathlib")} or \texttt{(load "mathlib.scm")},
the compiled version is used automatically---no source file is read,
no parsing occurs.

\section{Build System Integration}

The \texttt{schemecompiler} package provides a Quake template
(\texttt{scheme.tmpl}) with the following entry points:

\begin{center}
\begin{tabular}{lp{8cm}}
\toprule
\textbf{Directive} & \textbf{Effect} \\
\midrule
\texttt{scheme("mathlib")} &
  Compile \texttt{mathlib.scm} to \texttt{mathlib.i3}/\texttt{mathlib.m3}.
  The interface is \texttt{HIDDEN} (internal to the package). \\
\texttt{Scheme("mathlib")} &
  Same, but the interface is \texttt{VISIBLE} (exported).
  Other packages can \texttt{IMPORT mathlib} and call
  \texttt{mathlib.Install(interp)} directly. \\
\texttt{\_scheme\_compile("MathLib", "math-library.scm", HIDDEN)} &
  Explicit form for when the module name differs from the
  \texttt{.scm} filename. \\
\texttt{LinkCompiledSchemeAs("MyCompiled")} &
  Generate a synthetic module named \texttt{MyCompiled} that imports
  all compiled Scheme modules in this package, ensuring their
  \texttt{BEGIN} blocks execute at startup (triggering registry
  registration).  Must appear after all
  \texttt{scheme()}/\texttt{Scheme()} calls.
  Each package that compiles Scheme \emph{must} use a unique name
  to avoid duplicate-unit errors at link time. \\
\texttt{LinkCompiledScheme()} &
  Shorthand for \texttt{LinkCompiledSchemeAs("SchemeCompiled")}.
  Only suitable when a single package in the program compiles
  Scheme; prefer the \texttt{As} variant in general. \\
\bottomrule
\end{tabular}
\end{center}

\paragraph{Hyphenated filenames.}
\texttt{scheme("fold-constants")} automatically converts hyphens to
underscores in the Modula-3 module name, generating
\texttt{fold\_constants.i3}/\texttt{.m3}.  The compiled module
registers under both the underscored name and the original hyphenated
filename, so \texttt{(load "fold-constants.scm")} finds the compiled
version at runtime.

The \texttt{scheme}/\texttt{Scheme} naming follows the CM3
convention where a capitalized directive exports the interface
(cf.\ \texttt{Module}/\texttt{module},
\texttt{Interface}/\texttt{interface}).

\subsection{What the Build Does}

When \texttt{cm3} processes \texttt{scheme("mathlib")}, the Quake
template:

\begin{enumerate}
\item Checks if \texttt{mathlib.i3} is stale relative to \texttt{mathlib.scm}.
\item If so, writes a small bootstrap script that loads the compiler
  and calls \texttt{(compile-file "mathlib.scm" "mathlib")}.
\item Invokes the MScheme interpreter (\texttt{mscheme}) on that
  script.
\item Registers the derived \texttt{mathlib.i3} and \texttt{mathlib.m3}
  with the build system for compilation and linking.
\end{enumerate}

Incremental builds only re-run the Scheme compiler when the source
\texttt{.scm} file changes.

\paragraph{Linking the compiled modules.}
The \texttt{LinkCompiledSchemeAs} directive generates a linker module
(e.g., \texttt{MyCompiled.i3}/\texttt{.m3}) that \texttt{IMPORT}s
every compiled Scheme module.  However, Modula-3 only initializes
modules reachable from \texttt{Main}'s import graph.  You must add an
explicit import in your main module:

\begin{lstlisting}[language=Modula3]
(* Force compiled Scheme modules to be linked and registered *)
IMPORT MyCompiled; <*NOWARN*>
\end{lstlisting}

\noindent Without this \texttt{IMPORT}, the compiled modules are dead
code: they are linked into the binary but never initialized, and the
interpreter will load \texttt{.scm} files from disk as if no
compilation had occurred.

\paragraph{Bootstrap note.}
The \texttt{scheme()} template requires a pre-installed \texttt{mscheme}
binary to invoke the compiler.  If \texttt{mscheme} is not found (e.g.,
during a bootstrap build before the interpreter has been compiled),
the template silently becomes a no-op: the \texttt{.scm} files are
not compiled, and will be loaded by the interpreter at runtime as
usual.  This allows packages that use \texttt{scheme()} to build even
when the compiler toolchain is not yet available.

\section{Transparent Loading via the Registry}\label{sec:registry}

Each compiled module registers itself at Modula-3 module initialization
time (in the \texttt{BEGIN} block of the generated \texttt{.m3} file):

\begin{lstlisting}[language=Modula3]
BEGIN
  SchemeCompiledRegistry.Register("mathlib", Install);
  SchemeCompiledRegistry.Register("mathlib.scm", Install);
END mathlib.
\end{lstlisting}

The interpreter's \texttt{LoadFile} procedure checks this registry
before looking for a file on disk or in the bundled defaults.  If
the name matches a registered module, the compiled \texttt{Install}
procedure is called directly.

This means:
\begin{itemize}
\item \texttt{(load "mathlib.scm")} uses the compiled version.
\item \texttt{(require-modules "mathlib")} uses the compiled version.
\item No source file needs to be present at runtime.
\item The deduplication list in \texttt{require-modules} prevents
  double-installation, just as it prevents double-loading.
\end{itemize}

The \texttt{SchemeCompiledRegistry} interface:

\begin{lstlisting}[language=Modula3]
INTERFACE SchemeCompiledRegistry;
IMPORT Scheme;

TYPE Installer = PROCEDURE(interp : Scheme.T) RAISES { Scheme.E };

PROCEDURE Register(name : TEXT; installer : Installer);
PROCEDURE Unregister(name : TEXT);
PROCEDURE Lookup(name : TEXT) : Installer;  (* NIL if not found *)
PROCEDURE Disable();  (* force interpreter mode *)
PROCEDURE Enable();

END SchemeCompiledRegistry.
\end{lstlisting}

\noindent The \texttt{Disable}/\texttt{Enable} procedures are useful
during development: calling \texttt{Disable()} causes all
\texttt{Lookup} calls to return \texttt{NIL}, forcing the interpreter
to load source files from disk even when compiled versions are
registered.  This lets you edit a \texttt{.scm} file and
\texttt{(load ...)} it without rebuilding.

\paragraph{Environment variable.}
Setting the \texttt{MSCHEME\_INTERPRETED} environment variable
(to any value) disables the compiled registry at startup.  This is
equivalent to calling \texttt{Disable()} before any Scheme code runs:

\begin{lstlisting}[language=sh]
env MSCHEME_INTERPRETED=1 ./myprogram -scm
\end{lstlisting}

\noindent This is useful for debugging or verifying that interpreted
and compiled behavior are identical without recompiling.

If you use \texttt{Scheme("mathlib")} (visible interface), you can also
call \texttt{mathlib.Install(interp)} explicitly from Modula-3 code.
This is useful for modules that must be installed before
the first \texttt{(require-modules ...)} call.

\subsection{Reload Semantics}\label{sec:reload}

Compiled modules remain registered for the lifetime of the process.
Every call to \texttt{(load "mathlib.scm")} invokes the compiled
\texttt{Install} procedure, even if the module has already been loaded
once.  The \texttt{Install} procedure re-runs completely each time:
passthrough forms (such as \texttt{(define a 0)}) are re-evaluated,
and compiled procedures are re-installed into the environment.

This preserves the standard Scheme reload semantics.  For example:

\begin{lstlisting}
> (load "mathlib.scm")    ; a is bound to 0
> (set! a 1)              ; mutate a
> a
1
> (load "mathlib.scm")    ; reload: a is reset to 0
> a
0
\end{lstlisting}

\noindent The behavior is identical whether \texttt{mathlib.scm} is
compiled or interpreted: the module is re-executed from the top, all
definitions are re-established, and any mutations made since the
previous load are overwritten.

Note that \texttt{require-modules} maintains its own deduplication
list, so \texttt{(require-modules "mathlib")} loads the module at most
once per session regardless of how many times it is called.
Reloading via \texttt{load} bypasses this deduplication and always
re-executes the module.

\section{Dependencies and Preamble}

When a compiled module's source file contains top-level
\texttt{require-modules} or \texttt{load} forms, the compiler
extracts them and emits them at the start of the \texttt{Install}
procedure as \texttt{loadEvalText} calls:

\begin{lstlisting}[language=Modula3]
PROCEDURE Install(interp : Scheme.T) RAISES { Scheme.E } =
  VAR ...
  BEGIN
    (* Dependencies *)
    EVAL interp.loadEvalText("(require-modules \"basic-defs\")");
    (* forward declarations *)
    ...
  END Install;
\end{lstlisting}

If a dependency is itself a compiled module, the registry intercepts
the \texttt{require-modules} call and uses the compiled version
recursively.  The \texttt{require-modules} deduplication list
prevents infinite loops and redundant loading.

\section{Supported Constructs}

The Level~1 compiler handles the following Scheme constructs:

\begin{center}
\begin{tabular}{ll}
\toprule
\textbf{Category} & \textbf{Constructs} \\
\midrule
Control flow & \texttt{if}, \texttt{cond}, \texttt{case}, \texttt{and}, \texttt{or}, \texttt{not} \\
Binding      & \texttt{let}, \texttt{let*}, \texttt{letrec}, named \texttt{let} \\
Sequencing   & \texttt{begin} \\
Mutation     & \texttt{set!} \\
Iteration    & \texttt{do}, self-tail-calls (compiled to \texttt{LOOP}) \\
Closures     & \texttt{lambda}, internal \texttt{define} \\
Parameters   & Fixed arity, rest parameters (\texttt{.}~notation) \\
Constants    & Numbers, strings, booleans, characters, symbols, quoted lists \\
\bottomrule
\end{tabular}
\end{center}

Top-level \texttt{define} forms whose bodies use only
these constructs are compiled.  Forms using unsupported constructs
(\texttt{eval}, \texttt{current-environment}, reflective closure
access) are silently skipped---the interpreter handles them as usual
at load time.  The compiler reports which definitions were compiled
and which were skipped.

\section{Performance}

On ARM64\_DARWIN (Apple M4), with all optimizations enabled
(primitive inlining, direct calls, numeric unboxing):

\begin{center}
\begin{tabular}{lccc}
\toprule
\textbf{Benchmark} & \textbf{Compiled} & \textbf{Interpreted} & \textbf{Speedup} \\
\midrule
\texttt{(fibonacci 25)}       & 141\,ms    & 2,324\,ms   & 16.5$\times$ \\
\texttt{(ackermann 3 4)}      & 2.16\,ms   & 81\,ms      & 37.8$\times$ \\
\texttt{(loop-sum 10000 0)}   & 0.26\,ms   & 83\,ms      & 314$\times$ \\
\bottomrule
\end{tabular}
\end{center}

Self-tail-calls are converted to Modula-3 \texttt{LOOP} statements,
and numeric loop variables are unboxed to raw \texttt{LONGREAL},
eliminating all heap allocation in the loop body.  The 314$\times$
speedup for \texttt{loop-sum} reflects the combined effect of loop
conversion, primitive inlining, and numeric unboxing.

Compiled and interpreted code coexist freely.  A compiled function
can call an interpreted one (through the environment binding), and
vice versa.  This enables incremental compilation: compile the
hot spots and leave the rest interpreted.

\section{Limitations}

\begin{itemize}
\item \textbf{Top-level defines only.}  The compiler processes
  \texttt{(define (name params...) body)} forms at the top level.
  It does not compile standalone expressions, variable definitions,
  or macro definitions.

\item \textbf{No general tail-call optimization.}  Only self-tail-calls
  are optimized to \texttt{LOOP}.  Mutual recursion and tail calls to
  other functions use normal Modula-3 calls (consuming stack).

\item \textbf{No \texttt{eval} or \texttt{current-environment}.}
  Definitions that use these forms are silently left to the
  interpreter.

\item \textbf{No \texttt{call/cc}.}  MScheme's continuations are
  limited to stack unwinding; compiled code has the same limitation.

\item \textbf{Redefinition cost.}  Inter-module calls go through a
  binding indirection (\texttt{binding.get()} followed by
  \texttt{NARROW} and virtual dispatch).  For \emph{intra}-module
  calls, the compiler emits a fast-path pointer comparison: if the
  binding still holds the compiled procedure, a direct Modula-3
  call is made, bypassing virtual dispatch entirely.  This
  preserves correct redefinition semantics while delivering
  near-native call performance within a module.
\end{itemize}


%======================================================================
\chapter{Real-World Usage Patterns}
%======================================================================

\section{The Hybrid Application Pattern}\label{sec:hybrid}

The most effective way to use MScheme is as a hybrid architecture:
performance-critical libraries are coded in Modula-3 (or wrapped
from~C), while the application itself is assembled in Scheme.  The
Modula-3 code provides speed, type safety, and access to low-level
facilities; the Scheme code provides rapid iteration, exploratory
programming, and concise high-level composition.

This pattern is particularly powerful for scientific and engineering
applications, where numerically intensive kernels must be fast but
the top-level workflow changes frequently during development.

\subsection{Architecture}

A typical hybrid application has three components:

\begin{enumerate}
\item \textbf{M3 library modules} --- Compiled, type-safe,
  performance-critical code.  Numerical algorithms, data tables,
  I/O parsers, and C library wrappers.
\item \textbf{A thin M3 main program} --- Registers stubs, creates
  a Scheme interpreter, and enters the REPL or loads the application
  script.
\item \textbf{Scheme application code} --- Orchestrates the
  libraries, implements the top-level algorithm, handles
  command-line arguments, and produces output.
\end{enumerate}

\noindent The M3 code changes rarely and is compiled once.  The
Scheme code can be edited and reloaded interactively without
recompilation.

\subsection{Example: {\tt photopic}}

The {\tt photopic} application computes photometric quantities
for visible light: luminous efficacy, CIE color coordinates, Color
Rendering Index (CRI), and optimal light spectra via constrained
optimization.

\paragraph{M3 libraries (compiled, fast).}
The performance-critical code is in Modula-3 modules that provide
lookup tables with interpolation and physics computations:

\begin{center}
\begin{tabular}{lp{3in}}
\toprule
\textbf{Module} & \textbf{Role} \\
\midrule
{\tt Blackbody} & Planck radiance $B_\nu(T,\nu)$ \\
{\tt CieSpectrum} & CIE photopic/scotopic efficiency (41-point table, interpolated) \\
{\tt CieXyz} & CIE 1931 XYZ color matching functions (471-point table) \\
{\tt Tcs} & Test Color Sample reflectances (14 samples $\times$ 471 wavelengths) \\
{\tt TempUv} & Correlated Color Temperature $\leftrightarrow$ MacAdam~$uv$ lookup \\
{\tt ParametricSpectrum} & Parameterized spectrum representation for optimization \\
\bottomrule
\end{tabular}
\end{center}

\noindent These modules are pure Modula-3, compiled to native code.
The inner loops---integrating a spectrum against the CIE color
matching functions, for instance---run at full compiled speed.

\paragraph{The main program (thin M3 shell).}
{\tt Main.m3} is minimal: it registers stubs, creates a Scheme
interpreter, and enters the REPL:

\begin{lstlisting}[language=Modula3]
MODULE Main;
IMPORT SchemeStubs, SchemeM3, SchemeReadLine, ReadLine;

BEGIN
  SchemeStubs.RegisterStubs();
  WITH scm = NEW(SchemeM3.T).init(
    ARRAY OF TEXT { "require", "m3" }^) DO
    SchemeReadLine.MainLoop(
      NEW(ReadLine.Default).init(), scm)
  END
END Main.
\end{lstlisting}

\paragraph{The {\tt m3makefile} (wiring it all together).}
The build file compiles the M3 modules and generates Scheme stubs
for every interface the application wants to call from Scheme:

\begin{lstlisting}[language={}]
import ("mscheme")
import ("modula3scheme")
import ("sstubgen")
import ("integrate")
import ("cobyla")

Module ("CieSpectrum")
Module ("CieXyz")
Module ("Blackbody")
Module ("Tcs")
Module ("TempUv")
Module ("ParametricSpectrum")

SchemeStubs ("CieSpectrum")
SchemeStubs ("CieXyz")
SchemeStubs ("Blackbody")
SchemeStubs ("Tcs")
SchemeStubs ("TempUv")
SchemeStubs ("ParametricSpectrum")
SchemeStubs ("Math")
SchemeStubs ("Fmt")
SchemeStubs ("NR4p4")
SchemeStubs ("COBYLA_M3")
% ... plus I/O stubs: FileWr, Wr, Rd, etc.

ExportSchemeStubs ("photopic")
importSchemeStubs ()

implementation ("Main")
program ("photopic")
\end{lstlisting}

\noindent Note how both domain-specific modules ({\tt CieSpectrum},
{\tt Blackbody}) and standard library modules ({\tt Math}, {\tt Fmt},
{\tt NR4p4}) get stubs generated.  After compilation, all are callable
from Scheme.

\paragraph{The Scheme application (the actual program).}
The entire application logic lives in {\tt photopic.scm}
($\sim$800~lines of Scheme).  It uses the compiled M3 modules as
building blocks:

\begin{lstlisting}[language=Scheme]
(require-modules "display" "m3")

;; Physical constants
(define c 299792458.0)       ; speed of light
(define h 6.62607015e-34)    ; Planck's constant
(define l0 380e-9)           ; blue limit of vision
(define l1 770e-9)           ; red limit of vision

;; Blackbody radiance per wavelength -- calls compiled M3
(define (make-Bl T)
  (lambda (l)
    (let ((nu (/ c l)))
      (/ (Blackbody.PlanckRadiance T nu)
         (/ c (* nu nu))))))

;; Numerical integration -- calls compiled M3 (Romberg)
(define (integrate f a b)
  (NR4p4.QromoMidpoint (make-lrfunc-obj f) a b))

;; Luminous efficacy -- integrates spectrum against CIE data
(define (calc-lumens f)
  (integrate
    (lambda (x) (* (f x) (CieSpectrum.PhotoConv x)))
    l0 l1))

;; CIE XYZ color coordinates
(define (calc-Yxy spectrum)
  (let* ((xint (get-channel-integral spectrum 'x))
         (yint (get-channel-integral spectrum 'y))
         (zint (get-channel-integral spectrum 'z)))
    (xyz->Yxy (list xint yint zint))))
\end{lstlisting}

\noindent The CRI calculation---the core algorithm---is about
50~lines of Scheme that composes these primitives.  For each of
14~test color samples, it computes reflected spectra under reference
and test illuminants, transforms through CIE color spaces, and
computes color differences.  All the inner-loop numerical work
(interpolation, Planck's law, integration) runs as compiled M3; the
orchestration is Scheme.

\paragraph{Why this works well.}
\begin{itemize}
\item \textbf{Development speed:}\  The algorithm can be modified by
  editing {\tt photopic.scm} and calling {\tt (load "photopic.scm")}
  in the REPL---no recompilation, no relinking.
\item \textbf{Performance:}\  Tight numerical loops (spectrum
  integration, table interpolation) run at compiled speed.
  Only the top-level control flow pays the Scheme overhead.
\item \textbf{Exploratory use:}\  From the REPL, one can
  interactively compute the CRI of any spectrum, plot data to files,
  or run optimization sweeps---without writing a {\tt main} procedure.
\item \textbf{Code generation:}\  The Scheme code can even generate
  Modula-3 source.  In {\tt photopic}, a Scheme function computes
  a temperature-to-UV lookup table and writes it out as a Modula-3
  {\tt CONST} array, which is then compiled into the next build.
\end{itemize}

\bigskip
The rest of this chapter documents specific patterns from the CSP
compiler ({\tt cspc}), the most extensive user of MScheme, with over
80 {\tt SchemeStubs} declarations and approximately 40 Scheme source
files.

\section{Application Startup}

\begin{lstlisting}[language=Scheme]
; setup.scm -- loaded at startup
(require-modules "basic-defs" "m3" "hashtable" "set" "display")

; Extend load path to find application Scheme files
(define *m3utils* (Env.Get "M3UTILS"))
(set! **scheme-load-path**
      (cons (string-append *m3utils* "/csp/src")
            **scheme-load-path**))

; Load the compiler
(load "cspc.scm")
\end{lstlisting}

\section{AST Construction}

Building Modula-3 AST nodes from Scheme:

\begin{lstlisting}[language=Scheme]
(CspAst.SkipStmt)
(CspAst.AssignmentStmt lhs-expr rhs-expr)
(CspAst.IntegerExpr 42)
(CspAst.BooleanExpr #t)
(CspAst.IdentifierExpr "x")

;; Building sequences
(define seq (init-seq 'CspStatementSeq.T))
(seq 'addhi stmt1)
(seq 'addhi stmt2)
(CspAst.SequentialStmt (seq '*m3*))
\end{lstlisting}

\section{Round-Trip Between Representations}

A common pattern uses Scheme s-expressions as an intermediate
representation, converting to/from Modula-3 AST nodes:

\begin{lstlisting}[language=Scheme]
;; M3 -> Scheme: get s-expression form
(define lisp-form
  ((obj-method-wrap m3-obj 'CspSyntax.T) 'lisp))

;; Transform in Scheme
(define edited (simplify-stmt lisp-form))

;; Scheme -> M3: convert back
(define new-m3 (convert-stmt edited))
\end{lstlisting}

\section{Code Generation}

The CSP compiler generates complete {\tt m3makefile} entries and
Modula-3 source from Scheme, demonstrating meta-programming capabilities:

\begin{lstlisting}[language=Scheme]
(define wr (FileWr.Open "src/m3makefile"))
(Wr.PutText wr "import (\"mscheme\")\n")

(map (lambda (mod)
       (Wr.PutText wr
         (string-append "Smodule (\"" mod "\")\n")))
     module-names)

(map (lambda (sym)
       (Wr.PutText wr
         (string-append "SchemeStubs(\""
           (symbol->string sym) "\")\n")))
     *the-stubs*)

(Wr.PutText wr "ExportSchemeStubs (\"sim\")\n")
(Wr.PutText wr "importSchemeStubs ()\n")
(Wr.Close wr)
\end{lstlisting}


%======================================================================
\chapter{Environment and Debugging}
%======================================================================

\section{The Global Environment}

Every {\tt Scheme.T} has a global environment where top-level
definitions live.  The special variable {\tt *the-global-environment*}
refers to it.

\section{Error Handling}

Errors raise {\tt Scheme.E} with a text message.  In the REPL, a
traceback is printed:

\begin{lstlisting}[language={}]
> (car 42)
EXCEPTION! expected a pair, got: 42
(called from) eval (car 42)
\end{lstlisting}

Control tracebacks with {\tt (enable-tracebacks!)} and
{\tt (disable-tracebacks!)}.

Control Modula-3 runtime error mapping:
\begin{lstlisting}[language=Scheme]
(set-rt-error-mapping! #f)
;; M3 runtime errors crash instead of becoming Scheme.E
\end{lstlisting}

\section{Pickling (State Persistence)}

Scheme environments can be saved and restored:

\begin{lstlisting}[language=Scheme]
;; Save
(define wr (FileWr.Open "state.pkl"))
(dump-environment wr)
(wr-close wr)

;; Restore
(define rd (FileRd.Open "state.pkl"))
(load-environment! rd)
\end{lstlisting}


%======================================================================
\chapter{Concurrency Model}
%======================================================================

A single {\tt Scheme.T} instance is always single-threaded.  Its
methods must not be called concurrently from different threads.
However, the system supports:

\begin{itemize}
\item Multiple interpreters sharing a \emph{synchronized} global
  environment (using {\tt SchemeNavigatorEnvironment}).
\item Thread-safe ``navigable'' environments for cooperative
  interpreters.
\item The {\tt at-run} and {\tt clock-run} primitives, which create
  background Modula-3 threads that evaluate Scheme expressions.
\item Export of the global environment as a Network Object
  ({\tt netobj-export-global-environment}), enabling remote Scheme
  access over the network following the DEC~SRC Network Objects
  protocol~\cite{netobj}.
\end{itemize}


%======================================================================
\appendix
\chapter{Quick Reference Card}

\begin{lstlisting}[language=Scheme]
;; Module system
(require-modules "m3" "basic-defs" "display" "hashtable")

;; Load path
(set! **scheme-load-path**
      (cons "dir" **scheme-load-path**))

;; M3 procedure calls (after loading "m3")
(InterfaceName.ProcName arg1 arg2 ...)

;; M3 objects
(define obj (new-modula-object 'TypeName.T))
(define w (obj-method-wrap obj 'TypeName.T))
(w 'method-name arg1 arg2 ...)
(w '*m3*)                    ; raw M3 object

;; Hash tables
(define h (make-hash-table size hash-fn))
(h 'add-entry! key value)
(h 'retrieve key)
(h 'keys)

;; Display
(dis "output" var dnl)       ; print with newline

;; Type introspection
(rttype-typecode obj)
(list-modula-types)
(list-modula-type-ops tc)

;; Timing
(time expr)
(timenow)

;; I/O
(define wr (FileWr.Open "file"))
(Wr.PutText wr "text")
(Wr.Close wr)

;; Shell
(system "command")
\end{lstlisting}

\chapter{{\tt m3makefile} Cheat Sheet}

\begin{lstlisting}[language={}]
% Import mscheme packages
import ("mscheme")
import ("modula3scheme")
import ("sstubgen")

% Generate stubs for an interface
SchemeStubs ("MyInterface")

% Module + stubs combined
Smodule ("MyModule")

% Interface + stubs combined
Sinterface ("MyInterface")

% Sequence + stubs
Ssequence ("Key", "Value")

% Export stubs from this package
ExportSchemeStubs ("my-program-name")

% Import stubs from other packages
importSchemeStubs()

% Compile Scheme source to Modula-3
import ("schemecompiler")
scheme ("mylib")             % compiles mylib.scm, HIDDEN interface
Scheme ("mylib")             % compiles mylib.scm, VISIBLE interface
LinkCompiledSchemeAs ("MyCompiled")  % after all scheme() calls

% Build the program
implementation ("Main")
program ("my-program-name")

% In Main.m3, add:
%   IMPORT MyCompiled; <*NOWARN*>
\end{lstlisting}


%======================================================================
\newpage
{
\frenchspacing
\begin{thebibliography}{99}

\bibitem{r4rs} William Clinger and Jonathan Rees, eds.
Revised$^4$ Report on the Algorithmic Language Scheme.  1991.

\bibitem{r5rs} Richard Kelsey, William Clinger, and Jonathan Rees, eds.
Revised$^5$ Report on the Algorithmic Language Scheme.  1998.

\bibitem{sicp2} Hal Abelson and Gerald Jay Sussman with Julie Sussman.
{\it Structure and Interpretation of Computer Programs,} second edition.
Cambridge, Mass.:\ MIT Press, 1996.

\bibitem{spwm3} Greg~Nelson, ed.  {\it Systems Programming with Modula-3.}
Englewood Cliffs, N.J.:\ Prentice-Hall, 1991.

\bibitem{netobj} Andrew~Birrell, Greg~Nelson, Susan~Owicki, and Edward~Wobber.
{\it Network Objects,} digital SRC Research Report~115.
Palo Alto, Calif.:\ Digital Equipment Corporation, 1994 (revised 1995).

\bibitem{m3tk} Mick Jordan.  An extensible programming environment for Modula-3.
{\it Proceedings of the Fourth ACM SIGSOFT Symposium on Software
Development Environments,} Irvine, Calif., 1990.

\bibitem{jscheme} Peter Norvig.  JScheme.
{\tt https://norvig.com/jscheme.html}

\bibitem{scheme-initial-report} Gerald Jay Sussman and Guy L.~Steele.
Scheme: An Interpreter For Extended Lambda Calculus.
MIT~AI Memo AIM-349.  December 1975.

\end{thebibliography}
}

\end{document}
